\clearpage
\begin{proofbox}
\textbf{统一前提：}设 $k=\dfrac ba>0$，$A=(x_1,kx_1)$，
$B=(x_2,-kx_2)$，$O=(0,0)$。直线 $\ell$ 与两条渐近线各有
一个截点且 $A\ne B$。以下给出六种证明：行列式、初中共底面积比、
\textbf{横轴分块}、直角伸缩、点线距离与斜率差推广。

\medskip
\textbf{方法一：最快证明【行列式面积，考场两步】。}
\begin{center}% C088 PDF teaching figure; exact coordinates, equal x/y units.
\begin{tikzpicture}[x=1cm,y=1cm,>=stealth,font=\small]
\draw[rounded corners=5pt,black!25] (-1.6,-3.1) rectangle (9.2,4.5);
\node[font=\bfseries] at (3.8,3.85) {行列式给平行四边形面积，三角形占一半};
\coordinate (O) at (0,0); \coordinate (A) at (4,2.4);
\coordinate (B) at (2,-1.2); \coordinate (D) at (6,1.2);
\fill[teal!17] (O)--(A)--(B)--cycle;
\fill[orange!12] (A)--(D)--(B)--cycle;
\draw[->,black!45] (-0.8,0)--(7,0) node[right] {$x$};
\draw[->,black!45] (0,-1.8)--(0,2.8) node[above] {$y$};
\draw[blue!65!black,very thick] (O)--(A);
\draw[violet!75!black,very thick,dashed] (O)--(B);
\draw[black!65,thick] (A)--(D)--(B);
\draw[orange!85!black,very thick] (A)--(B);
\fill (O) circle (2pt) node[above left=4pt] {$O$};
\fill (A) circle (2pt) node[above=6pt] {$A(x_1,y_1)$};
\fill (B) circle (2pt) node[below=6pt] {$B(x_2,y_2)$};
\fill (D) circle (2pt) node[right=7pt] {$D=A+B$};
\node[teal!80!black] at (2,.3) {$S$};
\node[orange!85!black] at (4,1) {$S$};
\node at (3.8,-2.45) {$[OADB]=|x_1y_2-x_2y_1|=2S$};

\end{tikzpicture}
\end{center}


\textbf{第1步（立刻用渐近线写坐标）：}
\[
 y_1=kx_1,\qquad y_2=-kx_2.
\]
\textbf{第2步（带入以原点为顶点的面积公式）：}
\[
\begin{aligned}
S_{\triangle AOB}
&=\frac12\big|x_1y_2-x_2y_1\big|\\
&=\frac12\big|x_1(-kx_2)-x_2(kx_1)\big|\\
&=\boxed{k|x_1x_2|=\frac ba|x_1x_2|}.
\end{aligned}
\]
\textbf{适用信号：}已知两条直线的斜率与交点横坐标，
图中 $D=A+B$ 补成平行四边形，$AB$ 将它平分；求面积时\textbf{不必先求 $\ell$ 的方程}。

\medskip

\end{proofbox}
\clearpage
\begin{proofbox}
\textbf{方法二：最低门槛证明【初中共底面积比，无行列式、无伸缩】。}
\begin{center}% C088 PDF teaching figure; exact coordinates, equal x/y units.
\begin{tikzpicture}[x=1cm,y=1cm,>=stealth,font=\small]
\draw[rounded corners=5pt,black!25] (-1.2,-5.05) rectangle (11.1,4.3);
\node[font=\bfseries] at (4.95,3.7) {两次相似：横坐标比 $\to$ 斜边比 $\to$ 高之比};
\node at (1.5,2.7) {① 向 $x$ 轴投影，比较 $OB,OC$};
\node at (7.75,2.7) {② 共底 $OA$，比较两条高};
\begin{scope}
\coordinate (O) at (0,0); \coordinate (A) at (3,1.8); \coordinate (C) at (3,-1.8);
\coordinate (B) at (1.5,-.9); \coordinate (F) at (1.5,0); \coordinate (E) at (3,0);
\fill[orange!15] (O)--(A)--(C)--cycle;
\fill[teal!23] (O)--(B)--(F)--cycle;
\draw[->,black!45] (-.45,0)--(3.95,0) node[right] {$x$};
\draw[blue!65!black,thick] (O)--(A);
\draw[violet!75!black,dashed,thick] (O)--(C);
\draw[orange!85!black,very thick] (A)--(C);
\node[right=6pt] at (3,.65) {$2kx_1$};
\draw[teal!80!black,thick,dashed] (B)--(F);
\draw (1.32,0)--(1.32,-.18)--(1.5,-.18);
\draw (2.82,0)--(2.82,-.18)--(3,-.18);
\draw[black!65] (.55,0) arc[start angle=0,end angle=-30.96375653,radius=.55];
\fill (F) circle (1.3pt) node[above=4pt] {$F$};
\fill (E) circle (1.3pt) node[below right=3pt] {$E$};
\fill (B) circle (1.7pt) node[below left=4pt] {$B$};
\fill (O) circle (1.7pt) node[above left=3pt] {$O$};
\fill (A) circle (1.7pt) node[above=4pt] {$A$};
\fill (C) circle (1.7pt) node[right=8pt] {$C$};
\draw[<->,teal!80!black] (0,-2.35)--(1.5,-2.35) node[midway,fill=white,inner sep=3pt] {$x_2$};
\draw[<->] (0,-3.0)--(3,-3.0) node[midway,fill=white,inner sep=3pt] {$x_1$};
\node at (1.5,-3.6) {$\triangle OBF\sim\triangle OCE$};
\node at (1.5,-4.45) {$\dfrac{OB}{OC}=\dfrac{OF}{OE}=\dfrac{x_2}{x_1}$};
\end{scope}
\begin{scope}[shift={(6.2,0)}]
\coordinate (O) at (0,0); \coordinate (A) at (3,1.8);
\coordinate (B) at (1.5,-.9); \coordinate (C) at (3,-1.8);
\coordinate (D) at ({12/17},{36/85}); \coordinate (H) at ({24/17},{72/85});
\fill[orange!12] (O)--(A)--(C)--cycle;
\fill[teal!23] (O)--(A)--(B)--cycle;
\draw[blue!65!black,very thick] (O)--(A);
\draw[violet!75!black,thick,dashed] (O)--(C);
\draw[orange!70!black] (A)--(C);
\draw[teal!80!black,thick] (A)--(B);
\draw[red!65!black,dashed,thick] (B)--(D) (C)--(H);
\pic[draw,black!65,angle radius=3mm] {right angle=B--D--A};
\pic[draw,black!65,angle radius=3mm] {right angle=C--H--A};
\fill (D) circle (1.3pt) node[above left=3pt] {$D$};
\fill (H) circle (1.3pt) node[above left=3pt] {$H$};
\pic[draw,black!65,angle radius=5.5mm] {angle=B--O--A};
\path (B)--(D) node[midway,sloped,above=4pt,text=red!65!black] {$h_B$};
\node[red!65!black] at (2.7,-.45) {$h_C$};
\fill (O) circle (1.7pt) node[above left=3pt] {$O$};
\fill (A) circle (1.7pt) node[above=4pt] {$A$};
\fill (B) circle (1.7pt) node[below left=4pt] {$B$};
\fill (C) circle (1.7pt) node[right=8pt] {$C$};
\node at (1.55,-2.65) {$BD=h_B,\quad CH=h_C$};
\node at (1.55,-3.6) {$\triangle OBD\sim\triangle OCH$};
\node at (1.55,-4.45) {$\dfrac{h_B}{h_C}=\dfrac{BD}{CH}=\dfrac{OB}{OC}$};
\end{scope}

\end{tikzpicture}
\end{center}


先设 $x_1>0,x_2>0$；令 $C=(x_1,-kx_1)$。
\textbf{第1步（能直接量出的底和高）：}
$A=(x_1,kx_1)$，$C=(x_1,-kx_1)$，
故 $AC=2kx_1$，$O$ 到直线 $AC$ 的距离为 $x_1$，于是
\[
S_{\triangle OAC}=\frac12\times 2kx_1\times x_1=kx_1^2.
\]
\textbf{第2步（两次相似，再用共底面积比）：}
左图中，$F=(x_2,0),E=(x_1,0)$ 分别为 $B,C$ 在 $x$ 轴上的垂足。
$\angle OFB=\angle OEC=90^\circ$，且 $\angle BOF=\angle COE$，所以
\[
\triangle OBF\sim\triangle OCE
\quad\Longrightarrow\quad
\frac{OB}{OC}=\frac{OF}{OE}=\frac{x_2}{x_1}.
\]
右图中，$D,H$ 分别为 $B,C$ 到 $OA$ 所在直线的垂足。当 $k\ne1$ 时，
$\angle ODB=\angle OHC=90^\circ$，且 $\angle BOD=\angle COH$，所以
\[
\triangle OBD\sim\triangle OCH
\quad\Longrightarrow\quad
\frac{h_B}{h_C}=\frac{BD}{CH}=\frac{OB}{OC}.
\]
若 $k=1$，则 $OA\perp OC$，直接有 $h_B=OB,h_C=OC$。
$\triangle OAB,\triangle OAC$ 共底 $OA$，故
\[
\frac{S_{\triangle OAB}}{S_{\triangle OAC}}
=\frac{h_B}{h_C}=\frac{OB}{OC}=\frac{OF}{OE}=\frac{x_2}{x_1},
\qquad S_{\triangle OAB}=kx_1x_2.
\]
\textbf{第3步（符号推广）：}其他象限长度比改为 $OB/OC=|x_2/x_1|$，统一为
\[
\boxed{S_{\triangle OAB}=k|x_1x_2|=\frac ba|x_1x_2|}.
\]
\end{proofbox}
\clearpage
\begin{proofbox}
\textbf{方法三：横轴分块法【相似求截距，再把上下两块相加】。}
\begin{center}% C088: x-axis intercept, similar right triangles, and additive areas.
% Exact model: k=1/2, x1=6, x2=3, c=2*x1*x2/(x1+x2)=4.
\begin{tikzpicture}[x=1cm,y=1cm,>=stealth,font=\small]
\draw[rounded corners=5pt,black!25] (-1.0,-3.65) rectangle (11.8,5.3);
\node[font=\bfseries] at (5.4,4.65) {沿 $x$ 轴拆成两块：先求交点 $C$，再共底相加};
\node at (5.4,3.95) {$x_1>x_2>0,\quad B_1<C<A_1,\quad \tan\theta=k$};
\coordinate (O) at (0,0);
\coordinate (A) at (6,3); \coordinate (B) at (3,-1.5);
\coordinate (C) at (4,0);
\coordinate (A1) at (6,0); \coordinate (B1) at (3,0);
\fill[orange!17] (O)--(C)--(A)--cycle;
\fill[teal!23] (O)--(B)--(C)--cycle;
\draw[->,black!45] (-.55,0)--(10.85,0) node[right] {$x$};
\draw[->,black!45] (0,-1.65)--(0,3.2) node[above] {$y$};
\draw[blue!65!black,thick] (O)--(A);
\draw[violet!75!black,dashed,thick] (O)--(B);
\draw[teal!80!black,very thick] (2.8,-1.8)--(6.2,3.3);
\draw[red!65!black,very thick] (O)--(C);
\draw[red!65!black,dashed,thick] (A)--(A1) (B)--(B1);
\draw[black!65] (5.8,0)--(5.8,.2)--(6,.2);
\draw[black!65] (3.2,0)--(3.2,-.2)--(3,-.2);
\pic[draw,black!65,angle radius=8mm] {angle=C--O--A};
\pic[draw,black!65,angle radius=8mm] {angle=B--O--C};
\node at (1.2,.27) {$\theta$};
\node at (1.2,-.27) {$\theta$};
\fill (O) circle (1.7pt) node[below left=4pt] {$O$};
\fill (A) circle (1.7pt) node[above left=4pt] {$A$};
\fill (B) circle (1.7pt) node[right=5pt] {$B$};
\fill (C) circle (1.7pt) node[below=5pt] {$C$};
\fill (A1) circle (1.3pt) node[below=5pt] {$A_1$};
\fill (B1) circle (1.3pt) node[above=5pt] {$B_1$};
\path (A)--(A1) node[midway,right=5pt,text=red!65!black] {$kx_1$};
\path (B)--(B1) node[midway,left=5pt,text=red!65!black] {$kx_2$};
\node[orange!85!black] at (8.95,2.1) {$S_+=\frac12\,OC\cdot kx_1$};
\node[teal!80!black] at (8.95,-1.05) {$S_-=\frac12\,OC\cdot kx_2$};
\draw[<->] (0,-2.2)--(4,-2.2)
 node[midway,fill=white,inner sep=3pt] {$OC=c$};
\node at (5.4,-2.8) {$\triangle CA_1A\sim\triangle CB_1B$};
\end{tikzpicture}
\end{center}

\textbf{本方法重新记号：}$C=AB\cap x\text{轴}=(c,0)$，不沿用方法二的辅助点 $C$。
先取图中的 $x_1>x_2>0$；作 $AA_1\perp x$ 轴、$BB_1\perp x$ 轴，
其中 $A_1=(x_1,0),B_1=(x_2,0)$。令 $\theta$ 为渐近线与 $x$ 轴的锐角，则
$\tan\theta=k$，故 $AA_1=x_1\tan\theta=kx_1$、$BB_1=kx_2$。

\textbf{第1步（相似求 $C$）：}$A,B$ 分居 $x$ 轴两侧，$C$ 在线段 $AB$ 上，
且 $x_2<c<x_1$。两个直角三角形在 $C$ 处为对顶角，故
\[
\triangle CA_1A\sim\triangle CB_1B
\quad\Longrightarrow\quad
\frac{x_1-c}{c-x_2}=\frac{AA_1}{BB_1}=\frac{x_1}{x_2}.
\]
交叉相乘，得 $x_2(x_1-c)=x_1(c-x_2)$，所以
\[
\boxed{c=\frac{2x_1x_2}{x_1+x_2}},\qquad OC=c.
\]
\textbf{第2步（共底相加）：}$OC$ 将 $\triangle OAB$ 分成上方的 $\triangle OCA$
和下方的 $\triangle OCB$，两块共底 $OC$，于是
\[
S_{OAB}=\frac12\,OC\,(kx_1+kx_2)
=\frac12\cdot\frac{2x_1x_2}{x_1+x_2}\cdot k(x_1+x_2)
=\boxed{kx_1x_2}.
\]
\textbf{同号的其他位置：}$0<x_1<x_2$ 时交换左右位置，同理成立；
$x_1=x_2=t>0$ 时，$AB$ 竖直，直接有 $OC=t$、$S=\frac12t(2kt)=kt^2$。
两横坐标均负时，将全图绕 $O$ 旋转 $180^\circ$ 即归入上述情形。
\end{proofbox}
\clearpage
\begin{proofbox}
\textbf{方法三的适用边界：同号相加，异号相减，水平截线单独算。}
\begin{center}% C088: equal-scale boundary diagrams, k=1/2 in both panels.
\begin{tikzpicture}[x=.85cm,y=.85cm,>=stealth,font=\small]
\draw[rounded corners=5pt,black!25] (-.7,-2.85) rectangle (15.1,4.25);
\node[font=\bfseries] at (7.2,3.6) {横坐标异号：同侧作差；水平截线：直接算底高};
\node at (3,2.7) {$x_1=3,\ x_2=-1,\ c=-3$};
\node at (11,2.7) {$x_1=3,\ x_2=-3$};
\begin{scope}
\coordinate (C) at (0,0); \coordinate (O) at (3,0);
\coordinate (A) at (6,1.5); \coordinate (B) at (2,.5);
\fill[orange!17] (O)--(C)--(A)--cycle;
\fill[teal!23] (O)--(B)--(A)--cycle;
\draw[->,black!45] (-.35,0)--(6.35,0) node[right] {$x$};
\draw[blue!65!black,thick] (O)--(A);
\draw[violet!75!black,dashed,thick] (O)--(B);
\draw[teal!80!black,thick] (C)--(A);
\draw[red!65!black,dashed] (B)--(2,0) (A)--(6,0);
\draw[red!65!black,very thick] (C)--(O);
\fill (C) circle (1.7pt) node[below=5pt] {$C$};
\fill (O) circle (1.7pt) node[below=5pt] {$O$};
\fill (B) circle (1.7pt) node[above left=4pt] {$B$};
\fill (A) circle (1.7pt) node[above=4pt] {$A$};
\node at (3,-1.35) {$S_{OAB}=S_{OCA}-S_{OCB}$};
\end{scope}
\begin{scope}[shift={(8,0)}]
\coordinate (O) at (3,0); \coordinate (A) at (6,1.5); \coordinate (B) at (0,1.5);
\fill[teal!23] (O)--(A)--(B)--cycle;
\draw[->,black!45] (-.35,0)--(6.35,0) node[right] {$x$};
\draw[blue!65!black,thick] (O)--(A);
\draw[violet!75!black,dashed,thick] (O)--(B);
\draw[orange!85!black,very thick] (-.2,1.5)--(6.2,1.5);
\draw[red!65!black,dashed] (O)--(3,1.5) node[midway,right=4pt] {$k|x_1|$};
\draw[black!65] (3,1.3)--(3.2,1.3)--(3.2,1.5);
\fill (O) circle (1.7pt) node[below=5pt] {$O$};
\fill (A) circle (1.7pt) node[above=4pt] {$A$};
\fill (B) circle (1.7pt) node[above=4pt] {$B$};
\node at (3,-1.25) {$AB\parallel x\text{轴，无交点 }C$};
\node at (3,-2.05) {$S=\frac12\cdot2|x_1|\cdot k|x_1|$};
\end{scope}
\end{tikzpicture}
\end{center}

当 $x_1x_2<0$ 时，$A,B$ 在 $x$ 轴的\textbf{同一侧}，不能把两块面积相加。
若 $x_1+x_2\ne0$，$AB$ 所在直线仍与 $x$ 轴相交于
\[
C=(c,0),\qquad c=\frac{2x_1x_2}{x_1+x_2}.
\]
此时 $C$ 在线段 $AB$ 的延长线上；同样由竖直投影构成的直角三角形相似，
结合横轴上点的次序可得上式，也可由直线截距关系
\[
\frac{x_1-c}{kx_1}=\frac{x_2-c}{-kx_2}
\]
直接解出。$O,C,A,B$ 中两点 $O,C$ 在 $x$ 轴上，
$\triangle OCA,\triangle OCB$ 共底 $OC$ 且一大一小，故应作\textbf{面积之差}：
\[
\begin{aligned}
S_{OAB}
&=\bigl|S_{OCA}-S_{OCB}\bigr|
=\frac12|c|\,k\bigl||x_1|-|x_2|\bigr|\\
&=\frac12\left|\frac{2x_1x_2}{x_1+x_2}\right|k|x_1+x_2|
=\boxed{k|x_1x_2|}.
\end{aligned}
\]
若 $x_1+x_2=0$，则 $AB$ 是水平线，\textbf{不存在横轴交点 $C$}，不能使用 $OC$。
令 $x_1=t\ne0,x_2=-t$，直接取 $AB$ 为底：
\[
AB=2|t|,\qquad h=k|t|,\qquad
S_{OAB}=\frac12\cdot2|t|\cdot k|t|=kt^2=k|x_1x_2|.
\]
\textbf{选法提示：}横坐标同号时，沿横轴拆成上下两块最直观；
异号时若只求面积，可优先使用前面的行列式或共底面积比法。
\end{proofbox}
\clearpage
\begin{proofbox}
\textbf{方法四：直观补充【伸缩成直角】。}
\begin{center}% C088 PDF teaching figure; exact coordinates, equal x/y units.
\begin{tikzpicture}[x=0.9cm,y=0.9cm,>=stealth,font=\small]
\draw[rounded corners=5pt,black!25] (-1,-4.7) rectangle (12.6,5.6);
\node[font=\bfseries] at (5.8,4.95) {只伸缩纵坐标：两条渐近线变成直角};
\begin{scope}[shift={(0,0)}]
\node at (2,3.9) {原图：$k=2/3$};
\draw[->,black!45] (-0.5,0)--(4.3,0) node[right] {$x$};
\draw[->,black!45] (0,-3.45)--(0,2.6) node[above] {$y$};
\draw[blue!65!black,thick] (0,0)--(3.3,2.1999999999999997);
\draw[violet!75!black,thick,dashed] (0,0)--(3.3,-2.1999999999999997);
\coordinate (A) at (3,2.0); \coordinate (B) at (1.5,-1.0);
\fill[teal!18] (0,0)--(A)--(B)--cycle;
\draw[orange!85!black,thick] (A)--(B);
\draw[dotted,black!45] (A)--(3,0) (B)--(1.5,0);
\fill (0,0) circle (1.6pt) node[below left=4pt] {$O$};
\fill (A) circle (1.8pt) node[right=7pt] {$A$};
\fill (B) circle (1.8pt) node[below left=5pt] {$B$};
\node[below] at (3,0) {$x_1$}; \node[above] at (1.5,0) {$x_2$};
\node at (2,-4) {$S=3$};
\end{scope}
\begin{scope}[shift={(7,0)}]
\node at (2,3.9) {新图：$Y=y/k$，纵向放大 $3/2$};
\draw[->,black!45] (-0.5,0)--(4.3,0) node[right] {$x$};
\draw[->,black!45] (0,-3.45)--(0,2.6) node[above] {$Y$};
\draw[blue!65!black,thick] (0,0)--(3.3,3.3);
\draw[violet!75!black,thick,dashed] (0,0)--(3.3,-3.3);
\coordinate (A) at (3,3); \coordinate (B) at (1.5,-1.5);
\fill[teal!18] (0,0)--(A)--(B)--cycle;
\draw[orange!85!black,thick] (A)--(B);
\draw[dotted,black!45] (A)--(3,0) (B)--(1.5,0);
\fill (0,0) circle (1.6pt) node[below left=4pt] {$O$};
\fill (A) circle (1.8pt) node[right=7pt] {$A'$};
\fill (B) circle (1.8pt) node[below left=5pt] {$B'$};
\node[below] at (3,0) {$x_1$}; \node[above] at (1.5,0) {$x_2$};
\node at (2,-4) {$S'=9/2$，$S=kS'$};
\coordinate (A) at (3,3); \coordinate (B) at (1.5,-1.5); \coordinate (O) at (0,0);
\pic[draw,black!70,angle radius=4mm] {right angle=B--O--A};
\end{scope}

\end{tikzpicture}
\end{center}


\textbf{第1步：}纵向伸缩 $Y=y/k$，把 $y=\pm kx$ 变为 $Y=\pm x$。
于是 $A'=(x_1,x_1)$、$B'=(x_2,-x_2)$。

\textbf{第2步：}新图里 $OA'\perp OB'$，故
\[
S'=\frac12\cdot\underbrace{\sqrt2|x_1|}_{OA'}
\cdot\underbrace{\sqrt2|x_2|}_{OB'}=|x_1x_2|.
\]
\textbf{第3步：}新图的纵坐标与面积都变为原图的 $1/k$ 倍：
\[
S=kS'=\boxed{k|x_1x_2|}.
\]
\textbf{优势：}先看到直角三角形，再乘一个伸缩比例；
不存在“为什么 $\tfrac12$ 消失”的记忆负担。

\medskip
\end{proofbox}
\clearpage
\begin{proofbox}
\textbf{方法五：常规解析法【一边作底＋点到直线距离】。}
\begin{center}% C088 PDF teaching figure; exact coordinates, equal x/y units.
\begin{tikzpicture}[x=1cm,y=1cm,>=stealth,font=\small]
\draw[rounded corners=5pt,black!25] (-3,-4.05) rectangle (7.8,4.6);
\node[font=\bfseries] at (2.4,3.9499999999999997) {取 $OA$ 作底：垂足 $D$ 把点线距离变成高};
\coordinate (O) at (0,0); \coordinate (A) at (4,2.4); \coordinate (B) at (2,-1.2);
\coordinate (D) at ({16/17},{48/85});
\fill[teal!16] (O)--(A)--(B)--cycle;
\draw[->,black!45] (-1.8,0)--(5.5,0) node[right] {$x$};
\draw[->,black!45] (0,-.9)--(0,2.8) node[above] {$y$};
\draw[blue!65!black,thick] (-1.5,-0.8999999999999999)--(4.8,2.88);
\draw[violet!75!black,thick,dashed] (-1.5,.9)--(2.65,-1.59);
\draw[orange!85!black,very thick] (A)--(B);
\draw[blue!65!black,very thick] (O)--(A);
\draw[red!65!black,dashed,thick] (B)--(D);
\pic[draw,black!70,angle radius=3mm] {right angle=B--D--A};
\fill (O) circle (2pt);\node at (-.35,-.65) {$O$};
\fill (A) circle (2pt) node[above left=5pt] {$A$};
\fill (B) circle (2pt) node[below left=7pt] {$B$};
\fill (D) circle (1.7pt) node[above left=7pt] {$D$};
\node[red!65!black,left=8pt] at (1.47,-.32) {$h$};
\node[blue!65!black] at (-.75,-1.5) {$OA:\ kx-y=0$};
\node at (2.4,-2.25) {$OA=|x_1|\sqrt{1+k^2}$};
\node at (2.4,-3.15) {$BD=h=\dfrac{2k|x_2|}{\sqrt{1+k^2}}$};

\end{tikzpicture}
\end{center}


取 $OA$ 为底，所在直线 $kx-y=0$。由长度公式
\[
OA=\sqrt{x_1^2+k^2x_1^2}=|x_1|\sqrt{1+k^2}.
\]
由点到直线距离公式，$B(x_2,-kx_2)$ 到 $OA$ 的距离
\[
h=\frac{|kx_2-(-kx_2)|}{\sqrt{k^2+1}}
=\frac{2k|x_2|}{\sqrt{1+k^2}}.
\]
底乘高除以二：
\[
S_{\triangle AOB}=\frac12\,OA\cdot h
=\frac12\cdot|x_1|\sqrt{1+k^2}
\cdot\frac{2k|x_2|}{\sqrt{1+k^2}}
=\boxed{k|x_1x_2|}.
\]
\textbf{优势：}完全遵循高中解析几何的“底边长度＋点线距离”标准做法。

\medskip
\end{proofbox}
\clearpage
\begin{proofbox}
\textbf{方法六（补充）：双直线斜率差的统一面积模板。}
\begin{center}% C088 PDF teaching figure; exact coordinates, equal x/y units.
\begin{tikzpicture}[x=1cm,y=1cm,>=stealth,font=\small]
\draw[rounded corners=5pt,black!25] (-2.5,-2.6) rectangle (8.4,4.9);
\node[font=\bfseries] at (2.95,4.25) {两条过原点直线：竖直间隔体现斜率差};
\coordinate (O) at (0,0); \coordinate (A) at (3,3); \coordinate (B) at (2,.5);
\coordinate (Q) at (2,2);
\fill[teal!17] (O)--(A)--(B)--cycle;
\draw[->,black!45] (-1.1,0)--(5.4,0) node[right] {$x$};
\draw[->,black!45] (0,-0.7)--(0,3.1) node[above] {$y$};
\draw[blue!65!black,thick] (0,0)--(3.5,3.5);
\draw[violet!75!black,thick,dashed] (0,0)--(5,1.25);
\draw[orange!85!black,thick] (A)--(B);
\draw[red!65!black,very thick] (B)--(Q);
\node[red!65!black,left=6pt] at (2,1.25) {$d$};
\fill (O) circle (1.8pt) node[below left=4pt] {$O$};
\fill (A) circle (1.8pt) node[right=7pt] {$A$};
\fill (B) circle (1.8pt) node[below=20pt] {$B$};
\fill (Q) circle (1.8pt) node[above left=5pt] {$Q$};
\node[blue!65!black] at (-1.1,2.3) {$y=m_1x$};
\node[violet!75!black] at (4.8,.4) {$y=m_2x$};
\node at (2.95,-1.45) {$Q=(x_2,m_1x_2),\quad d=|m_1-m_2|\,|x_2|$};
\node at (2.95,-2.1) {图取 $m_1=1,\ m_2=1/4$；对称斜率差为 $2k$};

\end{tikzpicture}
\end{center}


若 $A=(x_1,m_1x_1)$、$B=(x_2,m_2x_2)$
分别在过原点的两条直线 $y=m_1x$、$y=m_2x$ 上，
其中 $m_1,m_2$ 有限且不同，$x_1x_2\ne0$。图中 $Q=(x_2,m_1x_2)$，
竖直间隔 $BQ=|m_1-m_2|\,|x_2|$，则
\[
\boxed{S_{\triangle AOB}=\frac{|m_1-m_2|}{2}|x_1x_2|}.
\]
证明仍只需两行：
\[
2S=|x_1m_2x_2-x_2m_1x_1|
=|m_2-m_1|\,|x_1x_2|.
\]
令 $m_1=k$、$m_2=-k$，由于 $|k-(-k)|=2k$，立刻得主结论。
\textbf{补充理由：}把双曲线的公式迁移为\textbf{任意两条过原点直线}的通解，
做斜率、面积类综合题时可直接复用。

\medskip
\textbf{必查边界：}水平或竖直截线也满足主公式；方法三须按前述边界分类处理。
若 $\ell$ 经过 $O$，它与两条渐近线只能同交于 $O$（除非重合，此时截点不唯一），
与题设的\textbf{两个不同截点}相矛盾。
\end{proofbox}
